% !TEX root = ../Thesis.tex

\chapter{Theoretische Fundierung}
\label{ch:theorie}

Moderne Softwarelösungen sind zunehmend durch verteilte Architekturen, Cloud-native Technologien und eine hohe Änderungsdynamik geprägt. Insbesondere Microservice-basierte Anwendungen mit Container-Orchestrierung und Continuous Deployment erhöhen die Komplexität des Betriebs erheblich \cite{sakinala2025observability}. Klassische Monitoring-Ansätze, die primär auf vordefinierte Metriken und statische Schwellwerte setzen, stoßen in solchen Umgebungen an ihre Grenzen, da sie häufig nur bekannte Fehlerzustände erfassen und wenig Unterstützung bei der Ursachenanalyse bieten \cite{sridharan2018distributed}.

Vor diesem Hintergrund hat sich der Ansatz der \emph{Observability} als wichtiges Konzept für den Betrieb moderner Systeme etabliert. Observability zielt darauf ab, den internen Zustand eines Systems anhand externer Messwerte nachvollziehbar zu machen und dadurch auch unbekannte oder unerwartete Fehlerszenarien analysierbar zu gestalten \cite{abieba2025advanced}.

\section{Begriffliche Grundlagen: Monitoring und Observability}
\label{sec:grundlagen}

Der Begriff Observability stammt ursprünglich aus der System- und Regelungstheorie und beschreibt die Fähigkeit, interne Zustände eines Systems aus dessen Ausgaben abzuleiten. Übertragen auf heutige Softwareanwendungen bezeichnet Observability die Eigenschaft, das Verhalten einer Application durch geeignete Telemetriedaten nachvollziehbar zu machen \cite{sridharan2018distributed}.

Monitoring und Observability sind dabei eng miteinander verbunden, jedoch nicht gleichzusetzen. Monitoring fokussiert sich auf die Überwachung bekannter Metriken und Ereignisse, um Abweichungen vom erwarteten Normalzustand zu erkennen. Observability erweitert diesen Ansatz, indem sie Kontextinformationen bereitstellt, die eine detaillierte Analyse von Ursachen und Zusammenhängen ermöglichen \cite{faseeha2025observability}. In der Literatur wird Observability daher häufig als Weiterentwicklung oder Übermenge des klassischen Monitorings beschrieben.

Zentral für Observability ist die systematische Erfassung unterschiedlicher Telemetriearten. Erst durch die kombinierte Betrachtung verschiedener Daten wird es möglich, komplexe Systemzustände sowie Ursachen von Fehlern zuverlässig zu analysieren \cite{abieba2025advanced,sakinala2025observability}.

\section{Kernbausteine der Observability}
\label{sec:kernbausteine}

In der Literatur werden vier Hauptelemente der Observability unterschieden: Metriken, Logs, Traces und Alerting. Diese Bausteine ergänzen sich gegenseitig und bilden gemeinsam die Grundlage für eine umfassende Analyse des Laufzeitverhaltens moderner Softwaresysteme \cite{abieba2025advanced,sakinala2025observability}.

\textbf{Metriken} sind numerische Messwerte, die den Zustand oder das Verhalten einer Anwendung über die Zeit hinweg beschreiben, beispielsweise Antwortzeiten, Fehlerraten oder Ressourcenauslastung. Sie eignen sich besonders für Trendanalysen, Kapazitätsplanung und die Definition von Service-Level-Zielen. Aufgrund ihrer aggregierten Form ermöglichen Metriken eine effiziente Überwachung auch großer Softwaresysteme \cite{sakinala2025observability}.

\textbf{Logs} bestehen aus ereignisbasierten Textinformationen, die detaillierte Einblicke in den Ablauf einzelner Systemkomponenten liefern. Sie enthalten kontextreiche Informationen, etwa Fehlermeldungen oder Zustandsänderungen, und sind insbesondere für die Ursachenanalyse und das Debugging grundlegend. Der Nachteil von Logs liegt in ihrem hohen Volumen und der vergleichsweise aufwändigen Auswertung \cite{abieba2025advanced}.

\textbf{Traces} dienen der Nachverfolgung einzelner Anfragen über mehrere Services hinweg. Sie zeigen, welche Komponenten an der Verarbeitung beteiligt waren und wie viel Zeit in den jeweiligen Verarbeitungsschritten verbracht wurde. Distributed Tracing ist insbesondere in Microservice-Architekturen ein essenzielles Mittel zur Identifikation von Latenzengpässen und Abhängigkeitsproblemen \cite{albuquerque2025tracing}.

\textbf{Alerting} ergänzt die passive Beobachtung durch aktive Benachrichtigung bei kritischen Zuständen. Alerts werden typischerweise auf Basis von Metriken oder Log-Ereignissen ausgelöst und sollen Betreiber frühzeitig auf Probleme hinweisen. Eine zentrale Herausforderung besteht darin, Alarmregeln so zu gestalten, dass relevante Vorfälle erkannt werden, ohne eine hohe Anzahl an Fehlalarmen zu erzeugen \cite{jalalvand2024alert}.

Erst das Zusammenspiel dieser vier Bausteine ermöglicht eine umfassende Observability. Während Metriken einen schnellen Überblick liefern, stellen Logs und Traces den notwendigen Detailgrad für die Ursachenanalyse bereit. Alerting sorgt schließlich dafür, dass Abweichungen zeitnah erkannt und adressiert werden können.

\section{Die Four Golden Signals nach Google SRE}
\label{sec:golden-signals}

Die \emph{Four Golden Signals} stellen ein technisches Referenzmodell zur Beobachtung von Diensten dar und wurden im Rahmen des Site Reliability Engineering von Google entwickelt \cite{ewaschukMonitoringDistributedSystems}. Sie beschreiben vier wesentliche technische Messgrößen, anhand derer sich Symptome von Leistungs- oder Stabilitätsproblemen identifizieren lassen:

\begin{itemize}
    \item \textbf{Latency}: Die Zeitspanne zur Bearbeitung einer Anfrage. Hohe oder stark schwankende Latenzen weisen auf Performanceprobleme hin.
    \item \textbf{Traffic}: Das Anfrage- oder Datenvolumen, das auf einen Dienst einwirkt, beispielsweise gemessen in Requests pro Sekunde.
    \item \textbf{Error}: Die Rate fehlgeschlagener Anfragen, etwa durch Serverfehler oder funktionale Fehlzustände.
    \item \textbf{Saturation}: Der Auslastungsgrad kritischer Ressourcen wie CPU, Speicher oder Warteschlangen.
\end{itemize}

Google empfiehlt, diese Signale als erste Orientierung im Monitoring zu nutzen, betont jedoch zugleich, dass ihre Aussagekraft stark vom jeweiligen Nutzungskontext abhängt \cite{ewaschukMonitoringDistributedSystems}. Insbesondere für Alerting und operative Entscheidungen ist eine Ergänzung durch Kontextinformationen erforderlich, etwa zur aktuellen Nutzungssituation, zur betroffenen Nutzergruppe oder zur zeitlichen Relevanz eines Problems. Ohne diese Einbettung besteht die Gefahr, dass technisch korrekte Anzeichen zu unnötigen oder nicht handlungsrelevanten Alarmierungen führen.

\section{Die Plattform frag.jetzt und ihre Systemarchitektur}
\label{sec:fragjetzt}

Die vorliegende Arbeit bezieht sich exemplarisch auf die Q\&A-Plattform \textit{frag.jetzt}, eine Cloud-native Webanwendung mit serviceorientierter Architektur. Die Plattform wurde an der Technischen Hochschule Mittelhessen entwickelt und dient der Durchführung interaktiver Fragerunden. Neben klassischer Frage-Antwort-Funktionalität integriert \textit{frag.jetzt} auch KI-gestützte Assistenzfunktionen zur automatisierten Beantwortung von Nutzeranfragen \cite{fragjetztRepo}.

Zur Einordnung der späteren Observability-Strategie wird im Folgenden ein Überblick über die Architekturkomponenten der Plattform gegeben \cite{fragjetztRepo}:

\begin{itemize}
    \item \textbf{Angular Frontend (SPA/PWA)}: Eine Single-Page-Application, über die Nutzer mit dem System interagieren. Das Frontend ist als Progressive Web App umgesetzt und kommuniziert überwiegend über HTTPS-APIs sowie WebSockets mit dem Backend.
    \item \textbf{WebSocket-Gateway (NGINX)}: Ein Gateway-Server zur Vermittlung der bidirektionalen Echtzeitkommunikation zwischen Frontend und Backend, insbesondere für das asynchrone Pushen von Antworten und Statusupdates.
    \item \textbf{Spring Boot Backend}: Der zentrale, Java-basierte Anwendungsserver, der die Kernlogik der Plattform implementiert. Er verarbeitet Nutzeranfragen, koordiniert interne Dienste und stellt REST- sowie WebSocket-Schnittstellen bereit.
    \item \textbf{Keycloak Identity Service}: Ein externer Dienst zur Authentifizierung und Autorisierung von Nutzern. Er übernimmt das Identitätsmanagement, führt jedoch zusätzliche Abhängigkeiten und potenzielle Latenzpunkte ein.
    \item \textbf{PostgreSQL-Datenbank}: Eine relationale Datenbank zur Persistierung des Anwendungszustands, etwa für Benutzerkonten, Fragen und Antworten. Aufgrund ihrer Rolle bei der Verarbeitung und Speicherung aller zustandsbehafteten Daten stellt sie einen kritischen Ausfallpunkt, einen sogenannten Single Point of Failure, bei hoher Last dar.
    \item \textbf{KI-Backend (Langchain \& FastAPI)}: Ein eigenständiger Microservice, der über Python (FastAPI) und die Langchain-Bibliothek Anfragen an externe Large-Language-Modelle weiterleitet, um KI-generierte Antworten zu erzeugen. Die Abhängigkeit von externen KI-APIs bringt variable Antwortzeiten und zusätzliche Fehlerquellen mit sich.
    \item \textbf{RabbitMQ Message Broker}: Ein Message-Broker zur asynchronen Entkopplung von Backend- und KI-Komponenten. RabbitMQ ermöglicht das Puffern und Weiterleiten von Antwortnachrichten und unterstützt so eine entkoppelte und skalierbare Verarbeitung.
\end{itemize}

Diese verteilten Komponenten bilden gemeinsam den durchgängigen Anfragepfad der Plattform. Eine Nutzerfrage durchläuft typischerweise mehrere Dienste, bevor eine Antwort an das Frontend zurückgegeben wird. Dadurch entstehen zahlreiche Schnittstellen und potenzielle Fehlerquellen. Eine Observability-Strategie für \textit{frag.jetzt} muss daher sicherstellen, dass entlang dieses gesamten Pfades geeignete Metriken, Logs und Traces erfasst werden, um Performanceprobleme und Ausfälle nachvollziehbar analysieren zu können \cite{faseeha2025observability}.

\section{Werkzeuge für Monitoring und Observability}
\label{sec:werkzeuge}

Zur Umsetzung von Monitoring- und Observability-Konzepten steht eine Vielzahl spezialisierter Open-Source-Werkzeuge zur Verfügung, die in modernen Cloud- und Microservice-Architekturen etabliert sind. Diese Werkzeuge adressieren unterschiedliche Aspekte der Telemetrieerfassung und -auswertung und werden häufig in Kombination eingesetzt \cite{banerjee2024comparative}.

\begin{itemize}
    \item \textbf{Prometheus}: Eine Open-Source-Software zur Erfassung und Speicherung zeitbasierter Metriken. Prometheus arbeitet mit einem Pull-basierten Ansatz, bei dem Metrik-Endpunkte regelmäßig abgefragt werden. Die integrierte Abfragesprache PromQL erlaubt detaillierte Analysen von Zeitreihendaten und bildet die Grundlage für regelbasierte Alerting-Mechanismen. Prometheus ist als CNCF-Projekt weit verbreitet und insbesondere für dynamische Cloud-Umgebungen geeignet \cite{banerjee2024comparative}.
    \item \textbf{Grafana}: Eine Plattform zur Visualisierung und Analyse von Observability-Daten aus unterschiedlichen Datenquellen. Grafana dient häufig als Dashboard-Werkzeug, in dem Metriken, Logs und Traces gemeinsam dargestellt und korreliert werden können. Dadurch unterstützt Grafana die visuelle Analyse von Systemzuständen und Trends \cite{giamattei2024monitoring}.
    \item \textbf{Elastic Stack (ELK)}: Der Elastic Stack kombiniert Elasticsearch, Logstash und Kibana und wird vor allem für die zentrale Aggregation und Analyse von Logdaten eingesetzt. Die Stärke des Stacks liegt in der skalierbaren Speicherung unstrukturierter Daten sowie in leistungsfähigen Such- und Analysefunktionen. Neben Logs unterstützt Elastic auch Metriken und Tracing, erfordert jedoch vergleichsweise hohe Systemressourcen \cite{banerjee2024comparative}.
    \item \textbf{Jaeger}: Ein verteiltes Tracing-Tool zur Nachverfolgung von Anfragen über Service-Grenzen hinweg. Jaeger ermöglicht die Analyse von End-to-End-Latenzen und unterstützt die Identifikation von Performanceengpässen sowie Abhängigkeitsbeziehungen in Microservice-Architekturen \cite{albuquerque2025tracing}.
    \item \textbf{OpenTelemetry}: Ein herstellerneutraler Standard zur Instrumentierung von Anwendungen. OpenTelemetry (OTel) stellt ein einheitliches API und SDK zur Erfassung von Metriken, Logs und Traces bereit und erlaubt die Weiterleitung dieser Daten an unterschiedliche Backends. Dadurch wird die Telemetrieerfassung von der Wahl konkreter Observability-Werkzeuge entkoppelt \cite{albuquerque2025tracing}.
\end{itemize}

Die Auswahl geeigneter Werkzeuge erfolgt typischerweise anhand mehrerer Bewertungskriterien, darunter die unterstützten Telemetriearten, Integrationsfähigkeit, Skalierbarkeit, Analyse- und Visualisierungsfunktionen sowie der Reifegrad der jeweiligen Lösung \cite{giamattei2024monitoring}. Diese Kriterien bilden die Grundlage für eine spätere vergleichende Analyse im Rahmen der Methodik.
